%% TOPTESI DEFAULT FRONTPAGE, IGNORED
%\ateneo{\thesisuniversity} % university name
%\logosede[5cm]{\thesisuniversitylogo} % logo, square brackets contain the height

%\titolo{\thesistitle} % title
%\sottotitolo{Metodo dei satelliti medicei} % subtitle

% place/remove a slash \\ to put the name on the following line or after Master Degree Course
%\corsodilaurea{\thesiscourse} % course name


%~251197 % id number is not needed

%\candidato{\thesiscandidatename~\textsc{\thesiscandidatesurname}} % candidate

% using tabular we can have more than 1 supervisor under the same column
%\relatore{\tabular{@{}l}%
%    \xmakefirstuc{\thesissupervisoronetitle}~\thesissupervisoronename~\textsc{\thesissupervisoronesurname}\\[0.4ex]
%    \xmakefirstuc{\thesissupervisortwotitle}~\thesissupervisortwoname~\textsc{\thesissupervisortwosurname}\\[0.4ex]
%    \xmakefirstuc{\thesissupervisorthreetitle}~\thesissupervisorthreename~\textsc{\thesissupervisorthreesurname}
%    \endtabular}
%\terzorelatore{Name}

% in this way we have Academic Year without stile=classica, so without lines
%\sedutadilaurea{\textsc{Academic~Year} 2019-2020}% per la laurea magistrale
% for PoliTo there is only month year
%\sedutadilaurea{\thesisdate}% per la laurea magistrale
% PhD
%\esamedidottorato{Novembre 1610}
%\ciclodidottorato{XV}

% offset for binding, the smaller the better
%\setbindingcorrection{3mm}


%\iflanguage{english}{%
	%\retrofrontespizio{This work is subject to the Creative Commons Licence}

%	\CorsoDiLaureaIn{\thesislevel's Degree Course in\space}

%	\TesiDiLaurea{\thesislevel's Degree Thesis}

%	\InName{in}
%	\CandidateName{\xmakefirstuc{\thesiscandidatetext}}% or Candidates
%	\AdvisorName{\xmakefirstuc{\thesissupervisortext}}% or Supervisor
	%\TutorName{Tutor}
	%\NomeTutoreAziendale{Internship Tutor}

	%\NomePrimoTomo{First volume}
	%\NomeSecondoTomo{Second Volume}
	%\NomeTerzoTomo{Third Volume}
	%\NomeQuartoTomo{Fourth Volume}
%}{}

\lstdefinelanguage{JavaScript}{
  keywords={break, case, catch, continue, debugger, default, delete, do, else, finally, for, function, if, in, instanceof, new, return, switch, this, throw, try, typeof, var, void, while, with, const, import, async, await},
  morecomment=[l]{//},
  morecomment=[s]{/*}{*/},
  morestring=[b]',
  morestring=[b]",
  sensitive=true
}

\lstset{
   language=JavaScript,
   extendedchars=true,
   basicstyle=\footnotesize\ttfamily,
   showstringspaces=false,
   showspaces=false,
   numbers=left,
   numberstyle=\footnotesize,
   numbersep=9pt,
   tabsize=2,
   breaklines=true,
   showtabs=false,
   captionpos=b
}

\lstdefinelanguage{JavaExtended}{
  language=Java, % Inherit base Java keywords (public, class, if, return, etc.)
  alsoletter={@}, % Tell LaTeX that @ is part of the keyword
  morekeywords={var, record, yield}, % Add modern Java keywords
  ndkeywords={ % Define your Spring, Quarkus, and Lombok annotations here
    @RestController, @RequestMapping, @GetMapping, @PostMapping, @PutMapping, @DeleteMapping,
    @Service, @Repository, @Component, @Autowired, @Bean, @Transactional, @CrossOrigin,
    @ModelAttribute, @RequestParam, @PathVariable, @RequestBody,
    @Path, @GET, @POST, @PUT, @DELETE, @ApplicationScoped, @RestPath, @RestQuery, @DefaultValue, @BeanParam,
    @Data, @Builder, @SuperBuilder, @NoArgsConstructor, @AllArgsConstructor, @RequiredArgsConstructor, @Slf4j,
    @MongoEntity, @Document, @Id, @Query,
    @Valid, @NotNull, @NotBlank, @Positive, @Pattern,
    @Override, @SuppressWarnings, @JsonIgnore, @JsonSerialize, @JsonDeserialize
  },
  % Define colors and styles
  keywordstyle=\color{blue}\bfseries,         % Standard Java keywords in blue
  ndkeywordstyle=\color{teal}\bfseries,       % Annotations in teal (or use purple/darkgreen)
  commentstyle=\color{gray}\itshape,          % Comments in gray italics
  stringstyle=\color{red}                     % Strings in red
}

\lstset{
   language=JavaExtended, % Use our newly defined language by default
   extendedchars=true,
   basicstyle=\footnotesize\ttfamily,
   showstringspaces=false,
   showspaces=false,
   numbers=left,
   numberstyle=\footnotesize\color{gray},
   numbersep=9pt,
   tabsize=2,
   breaklines=true,
   showtabs=false,
   captionpos=b,
   frame=single, % Adds a neat border around your code blocks
   rulecolor=\color{lightgray}
}